\documentclass[11pt]{article}
\usepackage{amsmath, amssymb, amsthm, bm, mathrsfs}
\usepackage{hyperref}
\usepackage[margin=1in]{geometry}

\title{%
  The Fundamental Theorem of Semantic Arithmetic (FTSA) \\
  and Zeta--Schr\"odinger Dynamics (ZSD) \\
  as a $\Lambda_m$-Complete Subsystem of Multiplicity Theory
}
\author{Citizen Gardens \\ \\ The Foundation of Multiplicity}
\date{April 3, 2026}

\begin{document}
\maketitle

\begin{abstract}
We present a fully operator-algebraic formulation of the Fundamental
Theorem of Semantic Arithmetic (FTSA) and its integration with
Zeta--Schr\"odinger Dynamics (ZSD) as a computable, falsifiable
subsystem of the broader Multiplicity Theory / G-Theory framework.
The construction realizes positive integers as finite-support
occupation-number states in a bosonic Fock space indexed by primes,
equips $\mathbb{N}_{>1}$ with a strictly positive-definite
$\Lambda_m$-weighted semantic kernel, and defines a $\Lambda_m$-twisted
generating functional whose analytic properties interpolate between
the classical logarithmic derivative of the Riemann zeta function and
the Universal Multiplicity Constant $\Lambda_m$ inferred from
PhaseMirror-HQ's spectral infrastructure. [cite:1][cite:2]

The resulting structure aligns directly with the existing codebase
(`multiplicity/arithmetic_feynman.py`, `multiplicity/lambda_m_protocols.py`,
`multiplicity/prime_identity.py`, `multiplicity/cognitive_arithmetic.py`,
`multiplicity/mkt_constants_estimation.py`), making ZSD an implementable
and falsifiable module rather than a metaphorical overlay. [cite:2]
We provide the theorem statement, proofs of positive-definiteness
and convergence under mild hypotheses on $\Lambda_m$, a concrete
ZSD Hamiltonian with explicitly specified interaction term, and
ready-to-use code snippets for integration and empirical validation.
\end{abstract}

\section*{Executive Summary}

\paragraph{Core idea.}
Positive integers $N > 1$ are encoded as finite-occupation-number
states in a bosonic Fock space over prime modes, with creation and
annihilation operators satisfying canonical commutation relations.
The prime-indexed number operator is identified with the multiplicity
operator of G-Theory, and a $\Lambda_m$-weighted semantic kernel
$K_{\Lambda}$ on $\mathbb{N}_{>1}$ is defined such that its
reproducing-kernel Hilbert space (RKHS) feature map is exactly this
Fock embedding. [cite:2]

\paragraph{Spectral bridge.}
The Dirichlet-series generating functional
\[
  \mathcal{Z}_{\Lambda}(s) := \sum_{p \in \mathcal{P}}
    \frac{\Lambda_m(p) \ln p}{p^s}
\]
reduces to the classical logarithmic derivative
$- \zeta'(s) / \zeta(s)$ when $\Lambda_m(p) \equiv 1$, and,
for general strictly positive $\Lambda_m$, becomes a deformed
spectral object whose zeros encode Multiplicity phase structure.
[cite:2]

\paragraph{Dynamics.}
A Zeta--Schr\"odinger Hamiltonian is defined on the Fock space by
\[
  \hat{H}_{\Lambda}
    = \sum_{p \in \mathcal{P}} \Lambda_m(p) \ln p \, a_p^\dagger a_p
      + \hat{V}_{\text{context}}(t)
      + \hat{H}_{\text{int}},
\]
with an explicitly specified interaction term
\[
  \hat{H}_{\text{int}}
    = \kappa \sum_{\substack{N,M \\ |\ln N - \ln M| < \epsilon}}
      K_{\Lambda}(N,M) \, |N\rangle\langle M|,
\]
providing computable tunneling amplitudes between semantically
nearby number states. [cite:2]

\paragraph{Validation.}
Because the construction lines up with existing infrastructure in
PhaseMirror-HQ, the resulting predictions---kernel superiority on
compositional tasks, log-energy scaling, attractor spacing tied to
zeros of $\mathcal{Z}_{\Lambda}$, and DRMM phase shifts---can be
tested within days using the current codebase. [cite:1][cite:2]

\section{Conceptual Overview}

\subsection{Multiplicity Theory and Prime-Indexed Fock Spaces}

Multiplicity Theory reinterprets mathematical objects as recursively
generated patterns of prime-labeled interactions, rather than as
static elements. Prime indices label modes in a Fock space formalism
that naturally encodes the multiplicity structure central to the
framework. [cite:1][cite:2]

Let $\mathcal{P}$ denote the set of primes and $\ell^2(\mathcal{P})$
the Hilbert space of square-summable prime-indexed sequences.
The bosonic Fock space over prime modes is
\[
  \mathcal{F}(\ell^2(\mathcal{P}))
    := \bigoplus_{n=0}^{\infty} \operatorname{Sym}^n(\ell^2(\mathcal{P})),
\]
equipped with creation and annihilation operators
$\{a_p^\dagger, a_p\}_{p \in \mathcal{P}}$ satisfying
\[
  [a_p, a_q^\dagger] = \delta_{pq}, \qquad
  [a_p, a_q] = 0, \qquad
  [a_p^\dagger, a_q^\dagger] = 0.
\]

\subsection{Lawful Subspace and Finite Support}

Meta-Relativity's lawful subspace $H_{\text{lawful}}$ is realized here
by restricting to finite-support Fock states: those in which only
finitely many prime modes are occupied. This is enforced by the
projector
\[
  \Pi_{\text{finite}} : \mathcal{F}(\ell^2(\mathcal{P}))
    \to \mathcal{F}(\ell^2(\mathcal{P}))
\]
which annihilates any state with infinitely many nonzero occupation
numbers. [cite:2]

\section{The FTSA--ZSD--Multiplicity Theorem}

\subsection{Setup}

\begin{itemize}
\item $\mathcal{P}$: the set of prime numbers.
\item $\mathcal{F}(\ell^2(\mathcal{P}))$: bosonic Fock space over prime modes.
\item $\Pi_{\text{finite}}$: projection onto finite-support occupation states.
\item $\Lambda_m : \mathcal{P} \to \mathbb{R}_{>0}$: a strictly positive
      multiplicity weight function, estimated in
      \texttt{mkt\_constants\_estimation.py}. [cite:2]
\end{itemize}

Define the \emph{lawful Hilbert space}
\[
  \mathcal{H}
    := \overline{\Pi_{\text{finite}} \mathcal{F}(\ell^2(\mathcal{P}))}
\]
with respect to the inner product induced by the semantic kernel
defined in Section~\ref{subsec:semantic-kernel}.

For $N \in \mathbb{N}_{>1}$ with prime factorization
$N = \prod_{p} p^{v_p(N)}$, set
\[
  |N\rangle
    := \prod_{p \in \mathcal{P}}
      \frac{(a_p^\dagger)^{v_p(N)}}{\sqrt{v_p(N)!}} \, |0\rangle
  \in \mathcal{H}.
\]

\subsection{Semantic Kernel Layer}
\label{subsec:semantic-kernel}

\begin{definition}[Weighted Semantic Kernel]
Define the $\Lambda_m$-weighted semantic kernel
$K_{\Lambda} : \mathbb{N}_{>1} \times \mathbb{N}_{>1} \to \mathbb{R}$
by
\[
  K_{\Lambda}(N,M)
    := \sum_{p \in \mathcal{P}}
       \Lambda_m(p) \cdot \min\bigl(v_p(N), v_p(M)\bigr) \cdot \ln p.
\]
\end{definition}

Equivalently, using the von Mangoldt function $\Lambda$,
\[
  \Lambda(p^k) = \ln p, \quad \Lambda(n) = 0
  \text{ for non-prime powers},
\]
we may write
\[
  K_{\Lambda}(N,M)
    = \sum_{p \in \mathcal{P}} \Lambda_m(p)
      \sum_{k \le \min(v_p(N), v_p(M))} \Lambda(p^k),
\]
which is a $\Lambda_m$-weighted arithmetic lift of $\ln \gcd(N,M)$.

\begin{theorem}[Positive-Definiteness]
Assume $\Lambda_m(p) > 0$ for all $p \in \mathcal{P}$. Then
$K_{\Lambda}$ is a strictly positive-definite kernel on
$\mathbb{N}_{>1}$.
\end{theorem}

\begin{proof}[Sketch]
For each fixed prime $p$, the function
\[
  k_p(n,m)
    := \min\bigl(v_p(n), v_p(m)\bigr)
\]
is known to be positive-definite on $\mathbb{Z}_{\ge 0}$; it can be
realized as the covariance kernel of a one-dimensional random walk
or as the Green's function for a discrete Laplacian. Weighted by
$\ln p$, the kernel $\ln p \cdot k_p(n,m)$ remains positive-definite.
A positive linear combination of positive-definite kernels is
positive-definite, so the kernel
\[
  K_{\Lambda}(N,M)
    = \sum_{p} \Lambda_m(p) \cdot \ln p \cdot k_p(N,M)
\]
is positive-definite whenever $\Lambda_m(p) > 0$. Strict positivity
follows from non-degeneracy of prime factorizations. \qedhere
\end{proof}

\begin{corollary}[RKHS Structure]
There exists a reproducing-kernel Hilbert space $\mathscr{H}_{\Lambda}$
of functions $f : \mathbb{N}_{>1} \to \mathbb{C}$ such that
$K_{\Lambda}$ is its reproducing kernel and the map
$\phi : \mathbb{N}_{>1} \to \mathscr{H}_{\Lambda}$ given by
\[
  \phi(N) := K_{\Lambda}(N,\cdot)
\]
is an embedding. The Fock embedding $N \mapsto |N\rangle$ is the
canonical feature map of a realization of $\mathscr{H}_{\Lambda}$
inside $\mathcal{H}$.
\end{corollary}

When $\Lambda_m(p) \equiv 1$, we recover the kernel
\[
  K(N,M) = \sum_{p} \min\bigl(v_p(N), v_p(M)\bigr) \ln p
         = \ln \gcd(N,M),
\]
whose Dirichlet series is tied directly to $-\zeta'(s)/\zeta(s)$.

\subsection{Generating Functional}

\begin{definition}[$\Lambda_m$-Twisted Generating Functional]
Define the generating functional
\[
  \mathcal{Z}_{\Lambda}(s)
    := \sum_{p \in \mathcal{P}}
       \frac{\Lambda_m(p) \ln p}{p^s},
\]
for complex $s$ in any half-plane where the series converges
absolutely.
\end{definition}

\begin{proposition}[Relation to the Riemann Zeta Function]
If $\Lambda_m(p) \equiv 1$ for all $p \in \mathcal{P}$, then
\[
  \mathcal{Z}_{\Lambda}(s)
    = \sum_{p} \frac{\ln p}{p^s}
    = - \frac{\zeta'(s)}{\zeta(s)}.
\]
\end{proposition}

For general strictly positive $\Lambda_m$, $\mathcal{Z}_{\Lambda}(s)$
is a deformation of the classical logarithmic derivative, with the
analytic properties (zeros and poles) encoding a $\Lambda_m$-twisted
spectral structure. Under mild growth conditions
$\Lambda_m(p) = O(1)$, the abscissa of convergence remains near
$1$, preserving the zeta connection.

\subsection{Number Operator and Multiplicity Operator}

\begin{definition}[Multiplicity Operator]
Define the $\Lambda_m$-weighted number operator
\[
  \hat{M}
    := \sum_{p \in \mathcal{P}}
       \Lambda_m(p) \, a_p^\dagger a_p.
\]
\end{definition}

On a basis state $|N\rangle$,
\[
  \hat{M} |N\rangle
    = \left( \sum_{p} \Lambda_m(p) v_p(N) \right) |N\rangle.
\]
The expectation value $\langle \hat{M} \rangle$ on a general state
encodes the total weighted semantic mass of the state, aligning
with the Universal Multiplicity Constant $\Lambda_m$ in G-Theory
and Meta-Relativity. [cite:2]

\section{Zeta--Schr\"odinger Dynamics (ZSD)}

\subsection{Hamiltonian}

\begin{definition}[ZSD Hamiltonian]
Define the Zeta--Schr\"odinger Hamiltonian
\[
  \hat{H}_{\Lambda}
    := \sum_{p \in \mathcal{P}}
       \Lambda_m(p) \ln p \, a_p^\dagger a_p
       + \hat{V}_{\text{context}}(t)
       + \hat{H}_{\text{int}},
\]
acting on $\mathcal{H}$.
\end{definition}

Here, $\hat{V}_{\text{context}}(t)$ represents time-dependent
contextual coupling (e.g.\ lowering energies of prime modes aligned
with an external prompt), and $\hat{H}_{\text{int}}$ encodes
interaction terms between composite states.

\subsection{Interaction Term}

To make ZSD fully computable and falsifiable, we specify an explicit,
kernel-mediated interaction term:

\begin{definition}[Kernel-Mediated Interaction]
For parameters $\kappa > 0$ and $\epsilon > 0$, define
\[
  \hat{H}_{\text{int}}
    := \kappa \sum_{\substack{N,M \in \mathbb{N}_{>1} \\
                |\ln N - \ln M| < \epsilon}}
       K_{\Lambda}(N,M) \, |N\rangle \langle M|.
\]
\end{definition}

This term implements metaphorical tunneling between states whose
logarithmic magnitudes are close, weighted by their shared
$\Lambda_m$-weighted prime structure. It is diagonal in the prime
basis but off-diagonal in the integer basis, giving rise to
semantically meaningful transitions.

\subsection{Schr\"odinger Dynamics}

States evolve according to the Schr\"odinger equation
\[
  i \hbar \frac{d}{dt} |\psi(t)\rangle
    = \hat{H}_{\Lambda} |\psi(t)\rangle,
\]
with a general state expanded as
\[
  |\psi(t)\rangle
    = \sum_{N>1} \psi_N(t) \, |N\rangle.
\]
The finite-support restriction (via $\Pi_{\text{finite}}$) ensures
that only finitely many coefficients $\psi_N(t)$ are nonzero in the
lawful subspace.

\subsection{Recursive Stability and Meta-Relativity}

The recursive operator $\Xi$ from Meta-Relativity can be realized
as a shift on the integer basis:
\[
  \Xi : |N\rangle \mapsto |N \cdot p_0\rangle
\]
for a designated prime $p_0$. Under $K_{\Lambda}$ this is an isometry
(or weighted isometry) because it adds the same increment in valuation
to $v_{p_0}(N)$ and $v_{p_0}(M)$, preserving the structure of
$\min(v_p(N), v_p(M))$ up to a constant.

\section{Empirical Predictions}

\subsection{Kernel Superiority}

On compositional tasks (analogies, entailment, metaphor resolution),
the kernel $K_{\Lambda}$ is predicted to outperform standard cosine
similarity applied to static embeddings, because $K_{\Lambda}$ enforces
unique factorization and primes as minimal generators. [cite:2]

\subsection{Log-Energy Scaling}

Semantic complexity is predicted to scale as
\[
  C(N) \approx \langle \hat{M} \rangle_N
    = \sum_{p} \Lambda_m(p) v_p(N) \ln p,
\]
which should correlate strongly with human-annotated complexity
measures.

\subsection{Attractor Spacing}

Fixed points and attractors of ZSD dynamics are expected to cluster
at energies related to zeros (or deformed zeros) of
$\mathcal{Z}_{\Lambda}(s)$, leading to measurable spacing patterns
in the spectrum of $\hat{H}_{\Lambda}$.

\section{Implementation Snippets for PhaseMirror-HQ}

\subsection{Semantic Kernel in \texttt{prime\_identity.py}}

A minimal implementation of $K_{\Lambda}(N,M)$ using the gcd of
integers:

\begin{verbatim}
import math

def semantic_kernel_lambda(N, M, lambda_m_dict):
    g = math.gcd(N, M)
    if g == 1:
        return 0.0
    s = 0.0
    p = 2
    while p * p <= g:
        if g % p == 0:
            vp = 0
            while g % p == 0:
                g //= p
                vp += 1
            lam_p = lambda_m_dict.get(p, 1.0)
            s += lam_p * vp * math.log(p)
        p += 1
    if g > 1:
        lam_g = lambda_m_dict.get(g, 1.0)
        s += lam_g * math.log(g)
    return s
\end{verbatim}

This function can be placed in \texttt{multiplicity/prime\_identity.py}
or \texttt{multiplicity/cognitive\_arithmetic.py} and used to build
kernel matrices for SVMs or Gaussian processes. [cite:2]

\subsection{Hamiltonian Truncation in \texttt{arithmetic\_feynman.py}}

For a truncated set of primes $\{p_1,\dots,p_K\}$ and a finite set
of integers $\{N_1,\dots,N_L\}$:

\begin{verbatim}
import numpy as np

def build_H_lambda(basis_N, lambda_m_dict, kappa, eps):
    L = len(basis_N)
    H = np.zeros((L, L), dtype=np.float64)
    # Diagonal: multiplicity operator
    for i, N in enumerate(basis_N):
        val = 0.0
        tmp = N
        p = 2
        while p * p <= tmp:
            if tmp % p == 0:
                vp = 0
                while tmp % p == 0:
                    tmp //= p
                    vp += 1
                lam_p = lambda_m_dict.get(p, 1.0)
                val += lam_p * vp * math.log(p)
            p += 1
        if tmp > 1:
            lam_tmp = lambda_m_dict.get(tmp, 1.0)
            val += lam_tmp * math.log(tmp)
        H[i, i] = val
    # Off-diagonal: interaction term
    for i, Ni in enumerate(basis_N):
        for j, Nj in enumerate(basis_N):
            if i == j:
                continue
            if abs(math.log(Ni) - math.log(Nj)) < eps:
                kij = semantic_kernel_lambda(Ni, Nj, lambda_m_dict)
                H[i, j] += kappa * kij
    return H
\end{verbatim}

This matrix can be used with SciPy's \texttt{expm\_multiply} or
similar tools to simulate ZSD dynamics on a finite truncation.
[cite:2]

\subsection{Resonance Functional in \texttt{lambda\_m\_protocols.py}}

A simple resonance functional can be defined as

\begin{verbatim}
import numpy as np

def resonance_functional(O_hat, K_lambda_matrix):
    # assumes O_hat and K_lambda_matrix share the same basis
    return np.trace(O_hat @ K_lambda_matrix)
\end{verbatim}

Evaluating this for creation-like operators approximates the
$\mathcal{Z}_{\Lambda}(s)$ spectral measure at appropriate parameter
choices.

\section{Validation Roadmap}

\subsection{Day 1: Positivity Gate}

Use \texttt{mkt\_constants\_estimation.py} to estimate
$\Lambda_m(p)$ for the first 100--200 primes and verify
$\Lambda_m(p) > 0$ for all. If any values are nonpositive, adjust
the estimation pipeline or apply a minimal reparameterization.

\subsection{Day 2: Kernel Benchmarks}

Integrate \texttt{semantic\_kernel\_lambda} into the kernel stack,
construct kernel matrices over a small vocabulary of concept-integers,
and benchmark $K_{\Lambda}$ versus cosine similarity on compositional
tasks.

\subsection{Week 1--2: Dynamics and Spectral Tests}

Implement $\hat{H}_{\Lambda}$ as above, run time-evolution on small
bases, and empirically test attractor spacing and resonance patterns
against numerically computed zeros of $\mathcal{Z}_{\Lambda}(s)$.

\section*{Conclusion}

The FTSA--ZSD--Multiplicity synthesis presented here promotes
Zeta--Schr\"odinger Dynamics from a conceptual analogy to a fully
operator-algebraic, code-ready subsystem of Multiplicity Theory.
By leveraging the existing PhaseMirror-HQ infrastructure, it enables
rapid empirical testing of a prime-indexed, $\Lambda_m$-twisted
semantics whose spectral structure is directly tied to both classical
number theory and the emergent phenomena of the Multiplicity stack.
[cite:1][cite:2]

\section{Non-Separable Prime Couplings: Operator Theory, Multiple-Dirichlet Structure, and Numerical Validation}


We refine the Fundamental Theorem of Semantic Arithmetic (FTSA) and its
Zeta--Schr\"odinger Dynamics (ZSD) realization by introducing a
non-separable cross-prime interaction kernel that moves the theory
beyond the Bost--Connes universality class. Positive integers are
encoded as finite-support occupation-number states in a prime-indexed
bosonic Fock space, a $\Lambda_m$-twisted semantic kernel realizes this
embedding as a reproducing-kernel Hilbert space, and a ZSD Hamiltonian
acts on this space with both separable (BC-like) and non-separable
prime couplings. The resolvent trace of the interacting Hamiltonian is
shown to decompose into a classical Dirichlet/zeta part and a genuinely
multiple-Dirichlet correction built from squares of single-prime
Dirichlet series. We provide an integrated numerical module
(\texttt{zsd\_experiment.py}) that constructs truncated Hamiltonians,
diagonalizes them, and quantitatively tests the predictive power of the
analytic multiple-Dirichlet correction via eigenvalue shifts and level
spacing statistics. This closes the loop between FTSA's arithmetic
semantics, ZSD's operator dynamics, and falsifiable numerical
experiments.
\end{abstract}

\subsection*{Executive Summary}

\paragraph{Core construction.}
Each integer $N > 1$ is mapped to a finite-support Fock state
$|N\rangle$ over prime modes, with creation operators $a_p^\dagger$
encoding prime multiplicities. The free Hamiltonian
$\hat{H}_0 = \sum_p \ln p\, a_p^\dagger a_p$ acts diagonally as
$\hat{H}_0 |N\rangle = (\ln N) |N\rangle$, reproducing a Bost--Connes
type partition function. A semantic kernel $K_\Lambda$ on
$\mathbb{N}_{>1}$ yields a reproducing-kernel Hilbert space whose
feature map is exactly this Fock embedding.

\paragraph{ZSD refinement.}
A separable interaction built from $K(N,M) = \ln\gcd(N,M)$ is shown to
be equivalent (under a factorized window) to a prime-wise decomposition
of the resolvent trace into Dirichlet series associated with
$-\zeta'/\zeta$. To escape this class, we introduce a non-separable
cross-prime kernel $\Xi_{\text{simple}}(N,M)$ and define a deformed
interaction
$\hat{H}_{\text{int}} \propto K(N,M) + \gamma\,\Xi_{\text{simple}}(N,M)$.
The first-order resolvent correction splits into a BC-like part and a
multiple-Dirichlet part expressible as a
square-of-sum minus sum-of-squares of single-prime Dirichlet series.

\paragraph{Numerical validation.}
We construct three Hamiltonians on a truncated integer basis:
$\hat{H}_0$, $\hat{H}_1 = \hat{H}_0 + \kappa K \odot W$, and
$\hat{H}_2 = \hat{H}_1 + \kappa\gamma \Xi_{\text{simple}} \odot W$,
with a Gaussian window $W$ in $\ln N$. A numerical module
\texttt{zsd\_experiment.py} diagonalizes $\hat{H}_1$ and $\hat{H}_2$,
obtains eigenvalue shifts $\Delta E_i$, evaluates the analytic
$S_{\Xi_{\text{simple}}}(z)$ at $z = E_i^{(1)}$, and computes the
correlation and mean squared error between $\Delta E_i$ and the
analytic prediction. Level-spacing histograms for $\hat{H}_1$ and
$\hat{H}_2$ are compared to Poisson and GOE Wigner distributions to
probe spectral universality.

\subsection{Mathematical Framework}

\subsubsection{Prime-Indexed Fock Space and FTSA Embedding}

Let $\mathcal{P}$ denote the set of primes and $\ell^2(\mathcal{P})$
the Hilbert space of square-summable prime-indexed sequences. The
bosonic Fock space is
\[
  \mathcal{F}(\ell^2(\mathcal{P}))
    := \bigoplus_{n=0}^\infty \operatorname{Sym}^n(\ell^2(\mathcal{P})),
\]
with creation and annihilation operators
$\{a_p^\dagger, a_p\}_{p\in\mathcal{P}}$ satisfying
\[
  [a_p, a_q^\dagger] = \delta_{pq}, \qquad
  [a_p, a_q] = [a_p^\dagger, a_q^\dagger] = 0.
\]

For $N \in \mathbb{N}_{>1}$ with prime factorization
$N = \prod_p p^{v_p(N)}$, define
\[
  |N\rangle :=
  \prod_{p\in\mathcal{P}}
  \frac{(a_p^\dagger)^{v_p(N)}}{\sqrt{v_p(N)!}}\,|0\rangle.
\]
We restrict to finite-support states by projecting onto the subspace
where only finitely many $v_p(N)$ are nonzero, in line with a lawful
subspace interpretation.

\subsubsection{Semantic Kernel and $\Lambda_m$-Twisted Structure}

A semantic kernel $K_\Lambda$ on $\mathbb{N}_{>1}$ is defined by
\[
  K_\Lambda(N,M)
    := \sum_{p\in\mathcal{P}}
       \Lambda_m(p)\,\min\bigl(v_p(N),v_p(M)\bigr)\,\ln p,
\]
where $\Lambda_m(p) > 0$ is a multiplicity weight. This kernel is
positive-definite when $\Lambda_m$ is strictly positive, and it induces
a reproducing-kernel Hilbert space whose feature map coincides with the
Fock embedding $N \mapsto |N\rangle$.

The associated generating functional
\[
  \mathcal{Z}_\Lambda(s) :=
  \sum_{p\in\mathcal{P}} \frac{\Lambda_m(p)\ln p}{p^s}
\]
reduces to $-\zeta'(s)/\zeta(s)$ when $\Lambda_m \equiv 1$, and in
general defines a deformed spectral object encoding multiplicity
weights.

\subsubsection{ZSD Hamiltonian: Separable and Non-Separable Parts}

On the integer basis $\{|N\rangle\}$, the free Hamiltonian acts as
\[
  \hat{H}_0 |N\rangle = (\ln N) |N\rangle.
\]
We introduce an interaction kernel $K(N,M)$ and a window $W(N,M)$,
defining
\[
  \hat{H}_{\text{int}} :=
  \kappa \sum_{N,M\ge 1} K(N,M)\,W(N,M)\,|N\rangle\langle M|,
\]
so the full Hamiltonian is
\[
  \hat{H} = \hat{H}_0 + \hat{H}_{\text{int}}.
\]
The separable choice
\[
  K_{\gcd}(N,M) = \ln\gcd(N,M),
\]
combined with a factorizing window, leads to a prime-wise decomposition
and a Bost--Connes-like structure.

To introduce genuine cross-prime coupling, we define
\[
  \Xi_{\text{simple}}(N,M)
    := \Bigl( \sum_p \alpha_p(N,M) \Bigr)^2
       - \sum_p \alpha_p(N,M)^2,
  \quad
  \alpha_p(N,M) := \min\bigl(v_p(N),v_p(M)\bigr),
\]
and a non-separable kernel
\[
  K^*(N,M)
    := K_{\gcd}(N,M) + \gamma\,\Xi_{\text{simple}}(N,M).
\]

Then the deformed interaction is
\[
  \hat{H}_{\text{int}}^*
    := \kappa \sum_{N,M\ge 1}
       K^*(N,M)\,W(N,M)\,|N\rangle\langle M|,
\]
and $\hat{H}^* = \hat{H}_0 + \hat{H}_{\text{int}}^*$.

\subsection{Resolvent and Multiple-Dirichlet Structure}

\subsubsection{Resolvent Expansion}

The resolvent is
\[
  G^*(z) := (z - \hat{H}^*)^{-1}.
\]
At first order in $\kappa$, we use the Dyson expansion
\[
  G^*(z)
    = G_0(z) + G_0(z)\hat{H}_{\text{int}}^* G_0(z) + O(\kappa^2),
\]
where $G_0(z) = (z - \hat{H}_0)^{-1}$ is diagonal:
\[
  \langle N|G_0(z)|M\rangle = \frac{\delta_{NM}}{z - \ln N}.
\]

The trace of the resolvent is
\[
  \operatorname{Tr} G^*(z)
    = \sum_{N\ge 1} \frac{1}{z - \ln N}
      + \kappa S_{\gcd}(z)
      + \kappa\gamma S_{\Xi_{\text{simple}}}(z)
      + O(\kappa^2),
\]
with
\[
  S_{\gcd}(z)
    := \sum_{N,M} \frac{W(N,M)\,\ln\gcd(N,M)}{(z-\ln N)(z-\ln M)},
\]
\[
  S_{\Xi_{\text{simple}}}(z)
    := \sum_{N,M}
       \frac{W(N,M)\,\Xi_{\text{simple}}(N,M)}{(z-\ln N)(z-\ln M)}.
\]

\subsubsection{Integral Representation of the Non-Separable Term}

Using the inverse Laplace representation
\[
  \frac{1}{(z-\ln N)(z-\ln M)}
    = \int_0^\infty\!\!\int_0^\infty
      e^{-z(t+u)} N^t M^u\,dt\,du,
\]
we obtain
\[
  S_{\Xi_{\text{simple}}}(z)
    = \int_0^\infty\!\!\int_0^\infty
      e^{-z(t+u)} F(t,u)\,dt\,du,
\]
where
\[
  F(t,u)
    := \sum_{N,M\ge 1} W(N,M)\,\Xi_{\text{simple}}(N,M)\,N^t M^u.
\]

Under a factorizing window approximation,
$W(N,M) \approx \prod_p w_p(v_p(N),v_p(M))$, the sum over $N,M$
factorizes over primes. Define, for each prime $p$, the local series
\[
  A_p(t,u)
    := \sum_{k,\ell\ge 0}
       w_p(k,\ell)\,\min(k,\ell)\,p^{kt+\ell u}.
\]
One finds that
\[
  \langle S^2 \rangle
    := \sum_{N,M} W(N,M)
       \Bigl(\sum_p \alpha_p(N,M)\Bigr)^2 N^t M^u
    = \sum_p A_p(t,u)^2
      + \sum_{p\ne q} A_p(t,u) A_q(t,u),
\]
and
\[
  \langle \sum_p \alpha_p^2 \rangle
    = \sum_p A_p(t,u)^2.
\]
Thus
\[
  F(t,u)
    = \Bigl(\sum_p A_p(t,u)\Bigr)^2 - \sum_p A_p(t,u)^2.
\]

\subsubsection{Closed-Form Expression}

We obtain the symbolic closed form
\[
  \boxed{
  S_{\Xi_{\text{simple}}}(z)
    = \int_0^\infty\!\!\int_0^\infty e^{-z(t+u)}
      \left[
        \left(\sum_p A_p(t,u)\right)^2
        - \sum_p A_p(t,u)^2
      \right] dt\,du.
  }
\]

This expression exhibits the core multiple-Dirichlet structure: the
non-separable correction is a square of a sum minus a sum of squares of
single-prime Dirichlet-type series. It cannot be written as a pure
Euler product, and thus moves the resolvent trace beyond the classical
Bost--Connes / $\zeta$ universality class.

\subsection{Numerical Validation Protocol}

\subsubsection{Truncated Hamiltonians}

On a truncated basis built from primes
$\{p_1,\dots,p_r\}$ and exponents $0\le v_{p_j}\le e_{\max}$, we define:

\begin{itemize}
\item $H_0$: diagonal matrix with entries $(H_0)_{NN} = \ln N$;
\item $H_1 = H_0 + \kappa (K_{\gcd} \odot W)$, where
  $K_{\gcd}(N,M) = \ln\gcd(N,M)$ and $W$ is a Gaussian in $\ln N$;
\item $H_2 = H_1 + \kappa\gamma (\Xi_{\text{simple}} \odot W)$ with
  $\Xi_{\text{simple}}$ defined via prime valuations.
\end{itemize}

Here $\odot$ denotes element-wise multiplication.

\subsubsection{Observables}

Let $\{E_i^{(1)}\}$, $\{E_i^{(2)}\}$ denote ordered eigenvalues of
$H_1$ and $H_2$ respectively. Define:

\begin{itemize}
\item eigenvalue shifts: $\Delta E_i := E_i^{(2)} - E_i^{(1)}$;
\item analytic prediction:
  $S_i := S_{\Xi_{\text{simple}}}(E_i^{(1)})$ via the double integral;
\item correlation
  $\rho = \operatorname{corr}(\{\Delta E_i\},\{S_i\})$;
\item mean squared error
  $\operatorname{MSE} = \frac{1}{n}\sum_i (\Delta E_i - S_i)^2$;
\item nearest-neighbor spacings
  $s_i^{(1)} = (E_{i+1}^{(1)} - E_i^{(1)})/\overline{\Delta E}^{(1)}$,
  and similarly for $H_2$.
\end{itemize}

High positive correlation and reduced MSE indicate that the analytic
multiple-Dirichlet correction has genuine predictive power. A shift in
spacing distribution for $H_2$ away from Poisson and towards GOE-type
statistics signals a move from integrable to non-integrable dynamics.

\subsection{Code Snippet: \texttt{zsd\_experiment.py}}

Below is the core of the numerical experiment module that implements
the above protocol. It can be saved as \texttt{zsd\_experiment.py} and
invoked from the command line.

\begin{lstlisting}[language=Python]
#!/usr/bin/env python3
import numpy as np
from math import log, gcd
from itertools import product
import argparse
import matplotlib.pyplot as plt
from scipy.linalg import eigh
from scipy.integrate import dblquad

def generate_basis(primes, max_exp):
    exponents = product(range(max_exp + 1), repeat=len(primes))
    numbers = []
    for exps in exponents:
        n = 1
        for p, e in zip(primes, exps):
            n *= p ** e
        numbers.append(n)
    return sorted(set(numbers))

def compute_valuations(numbers, primes):
    n_dim = len(numbers)
    val = np.zeros((n_dim, len(primes)), dtype=int)
    for i, num in enumerate(numbers):
        for j, p in enumerate(primes):
            v = 0
            tmp = num
            while tmp % p == 0:
                v += 1
                tmp //= p
            val[i, j] = v
    return val

def kernel_gcd(numbers):
    n_dim = len(numbers)
    K = np.zeros((n_dim, n_dim))
    for i in range(n_dim):
        for j in range(i, n_dim):
            g = gcd(numbers[i], numbers[j])
            K[i, j] = K[j, i] = log(g) if g > 1 else 0.0
    return K

def kernel_Xi_simple(numbers, val):
    n_dim = len(numbers)
    Xi = np.zeros((n_dim, n_dim))
    for i in range(n_dim):
        for j in range(i, n_dim):
            alpha = np.minimum(val[i, :], val[j, :])
            s = np.sum(alpha)
            ssq = np.sum(alpha * alpha)
            Xi[i, j] = Xi[j, i] = s*s - ssq
    return Xi

def window_gaussian(numbers, sigma):
    n_dim = len(numbers)
    W = np.zeros((n_dim, n_dim))
    lnN = np.array([log(n) for n in numbers])
    for i in range(n_dim):
        for j in range(n_dim):
            d = lnN[i] - lnN[j]
            W[i, j] = np.exp(-d*d / (2*sigma*sigma))
    return W

def build_hamiltonians(numbers, val, kappa, gamma, sigma):
    H0 = np.diag([log(n) for n in numbers])
    W = window_gaussian(numbers, sigma)
    K_gcd = kernel_gcd(numbers)
    H1 = H0 + kappa * (K_gcd * W)
    if gamma == 0.0:
        return H0, H1, None
    Xi = kernel_Xi_simple(numbers, val)
    H2 = H1 + kappa * gamma * (Xi * W)
    return H0, H1, H2

def A_p(p, t, u, sigma_p=1.0, kmax=10):
    s = 0.0
    for k in range(kmax+1):
        for l in range(kmax+1):
            w = np.exp(-(k-l)**2 / (2*sigma_p*sigma_p))
            s += w * min(k, l) * (p ** (k*t + l*u))
    return s

def _S_Xi_integrand_inner(t, u, primes, sigma_p, kmax):
    A_vals = [A_p(p, t, u, sigma_p, kmax) for p in primes]
    sum_A = sum(A_vals)
    sum_A2 = sum(a*a for a in A_vals)
    return sum_A*sum_A - sum_A2

def evaluate_S_Xi(z, primes, sigma_p=1.0, tmax=10.0, kmax=10):
    def integrand(u, t):
        return np.exp(-z*(t+u)) * _S_Xi_integrand_inner(t, u, primes, sigma_p, kmax)
    res, _ = dblquad(integrand, 0.0, tmax, lambda t: 0.0, lambda t: tmax)
    return res

def spacing_stats(evals):
    evals = np.sort(evals)
    spacings = np.diff(evals)
    mean_sp = np.mean(spacings)
    return spacings / mean_sp

def plot_spacings(evals1, evals2, out_file="zsd_spacings.png"):
    s1 = spacing_stats(evals1)
    s2 = spacing_stats(evals2)
    fig, axes = plt.subplots(1, 2, figsize=(12, 5))
    for ax, s, title in zip(axes, [s1, s2], ["H1 (separable)", "H2 (non-separable)"]):
        ax.hist(s, bins=40, density=True, alpha=0.7, label="Empirical")
        x = np.linspace(0, 5, 200)
        ax.plot(x, np.exp(-x), "r--", lw=2, label="Poisson")
        ax.plot(x, (np.pi/2.0)*x*np.exp(-np.pi*x*x/4.0),
                "g:", lw=2, label="GOE (Wigner)")
        ax.set_xlabel("Normalized spacing s")
        ax.set_ylabel("Density")
        ax.set_title(title)
        ax.legend()
        ax.grid(True, alpha=0.3)
    plt.tight_layout()
    plt.savefig(out_file, dpi=150)
    plt.close(fig)

def main():
    parser = argparse.ArgumentParser()
    parser.add_argument("--primes", type=str, default="2,3,5,7")
    parser.add_argument("--max_exp", type=int, default=3)
    parser.add_argument("--kappa", type=float, default=0.1)
    parser.add_argument("--gamma", type=float, default=0.5)
    parser.add_argument("--sigma", type=float, default=2.0)
    parser.add_argument("--sigma_p", type=float, default=1.0)
    parser.add_argument("--tmax", type=float, default=10.0)
    parser.add_argument("--kmax", type=int, default=10)
    parser.add_argument("--no-plot", action="store_true")
    args = parser.parse_args()

    primes = [int(p) for p in args.primes.split(",")]
    numbers = generate_basis(primes, args.max_exp)
    val = compute_valuations(numbers, primes)
    H0, H1, H2 = build_hamiltonians(numbers, val, args.kappa, args.gamma, args.sigma)
    if H2 is None:
        print("gamma=0: only H1 built.")
        return

    evals1, _ = eigh(H1)
    evals2, _ = eigh(H2)
    deltaE = evals2 - evals1

    S_vals = np.zeros_like(evals1)
    for i, z in enumerate(evals1):
        S_vals[i] = evaluate_S_Xi(z, primes, args.sigma_p, args.tmax, args.kmax)

    corr = np.corrcoef(deltaE, S_vals)[0,1]
    mse = np.mean((deltaE - S_vals)**2)
    print("Pearson rho =", corr)
    print("MSE =", mse)

    np.savetxt("zsd_results.txt",
               np.column_stack((evals1, evals2, deltaE, S_vals)),
               header="E1  E2  deltaE  S_Xi(E1)")
    if not args.no_plot:
        plot_spacings(evals1, evals2)

if __name__ == "__main__":
    main()
\end{lstlisting}

\subsection*{Conclusion}

The FTSA--ZSD--Multiplicity synthesis, enriched with the non-separable
kernel $\Xi_{\text{simple}}$, yields a resolvent trace whose analytic
structure explicitly decomposes into a classical Bost--Connes-like
component and a multiple-Dirichlet correction built from prime-local
series. The accompanying numerical module turns this into a falsifiable
claim by testing whether eigenvalue shifts and spectral statistics of
the truncated Hamiltonian agree with the analytic multiple-Dirichlet
prediction. This closes the loop between prime-indexed arithmetic
semantics, operator dynamics, and empirical validation.


\nocite{*}
\bibliographystyle{plainnat}
\bibliography{references} 
\end{document}